\chapter{Product Requirements Document (PRD)}
\label{ch:prd}

\section{Product vision}

A developer should never hand-write a login screen again. Expo Snips aims to make
any common mobile screen available as a versioned, stack-aware recipe that can be
generated into an existing project in one action and kept current thereafter.

\textbf{Status:} INFERRED (vision statement); the underlying capability is
IMPLEMENTED.

\section{Problem statement}

See Section~2.3 of Chapter~\ref{ch:project-description}. In short: recurring
screen boilerplate is slow, inconsistent and easy to get subtly wrong, and no
existing category of tool (snippet library, project starter, component library)
generates whole, wired screens into an existing stack.

\section{Goals and objectives}

\begin{longtable}{p{6cm}p{8.6cm}}
\caption{Goals and how they are served}\\
\toprule
\textbf{Goal} & \textbf{Mechanism} \\
\midrule

Generate complete screens & Multi-file recipes (\texttt{files[]} with conditions). \\
Match the user's stack & \texttt{detectProject()} plus conditional template rendering. \\
Make output trustworthy & Best-practice libraries (react-hook-form, zod) and standard platform patterns. \\
Remove dependency friction & Package-manager detection and installation. \\
Support long-term maintenance & Manifest, \texttt{status}, diff-previewed \texttt{update}. \\
Work in terminal and editor & One engine, two surfaces. \\
Zero-config, offline-capable & Catalog bundled into both clients. \\
\bottomrule
\end{longtable}

\textbf{Status:} IMPLEMENTED (mechanisms) / RECOMMENDED (goal phrasing).

\section{Personas}

See Section~2.6 of Chapter~\ref{ch:project-description} for the three personas
(solo developer, team/agency developer, VS Code power user).

\section{User stories}

Stories are written in the standard form and each is mapped to the code that
satisfies it. Stories without an implementation are listed separately as future
scope.

\begin{longtable}{p{1.3cm}p{6.4cm}p{6.9cm}}
\caption{Implemented user stories}\\
\toprule
\textbf{ID} & \textbf{Story} & \textbf{Implementation} \\
\midrule

US-01 & As a developer, I want to add a screen with one command so I do not hand-write boilerplate. & \texttt{add.ts}: \texttt{expo-snips add <screenId>}. \\
US-02 & As a developer, I want the generated code to match my project's styling. & \texttt{detectProject.ts} + \texttt{\#if (eq styling ...)} in templates. \\
US-03 & As a developer, I want to choose a screen interactively if I do not know its id. & \texttt{add.ts} interactive \texttt{prompts} select when no argument is given. \\
US-04 & As a developer, I want form validation in generated auth screens. & \texttt{zod} schema + \texttt{zodResolver} in \texttt{components.login-form.hbs}. \\
US-05 & As a developer, I want loading and error states. & \texttt{isLoading} state and error rendering in the login template. \\
US-06 & As a developer, I want keyboard handling. & \texttt{KeyboardAvoidingView} with platform-specific behaviour. \\
US-07 & As a developer, I want the right dependencies installed automatically. & \texttt{installDeps.ts} / \texttt{generator.ts}. \\
US-08 & As a developer, I want to see which screens my project uses. & \texttt{expo-snips status} reading \texttt{.expo-snips.json}. \\
US-09 & As a developer, I want to update generated screens safely. & \texttt{update.ts} with \texttt{printDiff} and confirmation. \\
US-10 & As a VS Code user, I want to add a screen without leaving the editor. & \texttt{completionProvider.ts}, \texttt{commands.ts}, status bar in \texttt{extension.ts}. \\
US-11 & As a TypeScript user, I want TypeScript output. & Recipe \texttt{typescript} flag and \texttt{\#if typescript} blocks; TS/JS file variants. \\
US-12 & As a JavaScript user, I want JavaScript output. & \texttt{--javascript} flag; \texttt{hooks/useLogin.js} variant. \\
US-13 & As a developer using shadcn-style UI, I want the primitives generated. & \texttt{primitives.ts} plus \texttt{ui/*} recipes. \\
US-14 & As a developer, I want to start a brand-new Expo project. & \texttt{init.ts} running \texttt{create-expo-app}. \\
US-15 & As a maintainer, I want to validate recipes before publishing. & \texttt{scripts/validate.mjs} against \texttt{schema.json}. \\
\bottomrule
\end{longtable}

\section{Use cases}

\begin{longtable}{p{2.2cm}p{5.2cm}p{7.2cm}}
\caption{Use cases}\\
\toprule
\textbf{ID} & \textbf{Use case} & \textbf{Flow} \\
\midrule

UC-01 & Add an authentication screen & \texttt{add auth/login} $\rightarrow$ detect $\rightarrow$ render $\rightarrow$ install $\rightarrow$ manifest. \\
UC-02 & Add an app-chrome screen & \texttt{add app/home}. \\
UC-03 & Generate a glassmorphism variant & \texttt{add app/home --variant glass} (variant templates exist for \texttt{app/home}, \texttt{auth/login}, \texttt{onboarding/welcome}). \\
UC-04 & Generate shadcn UI & \texttt{add auth/register --ui shadcn} (forces NativeWind and lucide; adds primitives). \\
UC-05 & Generate JavaScript output & \texttt{add auth/login --javascript}. \\
UC-06 & Generate through react-navigation & \texttt{add auth/login --router react-navigation} (supported only for \texttt{auth/login} and \texttt{auth/register}). \\
UC-07 & Inspect installed screens & \texttt{status}. \\
UC-08 & Update installed screens & \texttt{update} or \texttt{update <id> --yes}. \\
UC-09 & Use a custom registry & \texttt{add app/home --registry http://localhost:3000}. \\
UC-10 & Bootstrap a project & \texttt{init --name myapp}. \\
\bottomrule
\end{longtable}

\section{Features}

The reference list of features is the recipe and command inventory in
Section~3.10 of Chapter~\ref{ch:project-description}. The catalog itself is
tabulated in Appendix~A.

\section{Functional requirements}

Functional requirements are numbered and detailed in
Chapter~\ref{ch:functional-requirements}. In the PRD they are summarised as:

\begin{itemize}
  \item FR-1 Catalog resolution (bundled or remote).
  \item FR-2 Project stack detection.
  \item FR-3 Option resolution and validation.
  \item FR-4 Conditional template rendering.
  \item FR-5 File writing.
  \item FR-6 UI primitive resolution.
  \item FR-7 Dependency installation.
  \item FR-8 Manifest maintenance.
  \item FR-9 Update with diff preview.
  \item FR-10 Interactive selection.
  \item FR-11 Project bootstrap.
\end{itemize}

\section{Non-functional requirements}

See Chapter~\ref{ch:non-functional-requirements}. Highlights: offline operation,
no telemetry, cross-platform paths, deterministic output, and a 10-second network
timeout.

\section{Acceptance criteria}

Acceptance criteria are stated against observable behaviour.

\begin{longtable}{p{1.3cm}p{12.6cm}}
\caption{Acceptance criteria}\\
\toprule
\textbf{ID} & \textbf{Criterion} \\
\midrule

AC-01 & Running \texttt{expo-snips add auth/login} in a TypeScript Expo Router project creates \texttt{app/login/index.tsx}, \texttt{components/LoginForm.tsx} and \texttt{hooks/useLogin.ts}. \\
AC-02 & The generated screen contains a \texttt{KeyboardAvoidingView} and a \texttt{SafeAreaView}. \\
AC-03 & Validation errors render for empty or malformed input because a \texttt{zod} schema is applied through \texttt{react-hook-form}. \\
AC-04 & With \texttt{--ui shadcn}, \texttt{components/ui/button.tsx}, \texttt{components/ui/input.tsx} and \texttt{lib/utils.ts} are created unless they already exist. \\
AC-05 & With \texttt{--no-install}, no package-manager command is executed. \\
AC-06 & After generation, \texttt{.expo-snips.json} contains an install entry for the recipe with its version and options. \\
AC-07 & \texttt{expo-snips status} lists every entry in the manifest. \\
AC-08 & \texttt{expo-snips update} prints a diff and changes nothing when the user declines. \\
AC-09 & In VS Code, typing \texttt{snip:} offers the recipe identifiers and selecting one generates the screen. \\
AC-10 & With no network and no registry configured, \texttt{list} still reports the bundled catalog. \\
AC-11 & \texttt{expo-snips add auth/login --styling tamagui} exits with an error because the recipe does not support it. \\
\bottomrule
\end{longtable}

\textbf{Status:} IMPLEMENTED (each criterion is directly derivable from the cited
source) --- except AC-11, which is the intended validation path in
\texttt{add.ts} lines 104--109.

\section{Success metrics}

\textbf{UNKNOWN.} See Section~4.9 of Chapter~\ref{ch:brd}. No metric data exists
in the repository.

\section{Constraints and dependencies}

\begin{itemize}
  \item Runtime: Node.js $\geq$ 18.
  \item Engine dependency: \texttt{handlebars} \^{}4.7.8 (only runtime dependency of
        \texttt{@expo-snips/core}).
  \item CLI dependencies: \texttt{commander} \^{}12, \texttt{prompts} \^{}2.4,
        \texttt{chalk} \^{}4.1, \texttt{fs-extra} \^{}11.2.
  \item Extension build: \texttt{esbuild} \^{}0.21.5 and \texttt{@vscode/vsce} \^{}2.24
        (development only).
  \item Generated screens depend on recipe-declared packages; for auth recipes
        that is \texttt{react-hook-form}, \texttt{zod}, \texttt{@hookform/resolvers}
        and \texttt{lucide-react-native}.
\end{itemize}

\section{Risks}

See Section~4.11 of Chapter~\ref{ch:brd} (R-01 \dots R-06).

\section{Future scope}

The following are \emph{not} implemented and are listed as candidate future work.
They are drawn from gaps in the repository and from the extension's own
``Roadmap'' section, which is itself an unimplemented checklist.

\begin{itemize}
  \item VS Code commands for update and status (the CLI has them; the extension
        does not).
  \item Tamagui styling support, Vector Icons support, custom UI primitive
        support (declared in types/schema/settings; no recipes).
  \item Test infrastructure (\texttt{README.md} ``Roadmap'': ``Testing
        infrastructure'').
  \item CI for packaging the VS Code extension (\texttt{README.md} ``Roadmap'':
        ``CI/CD pipeline for the VS Code extension'').
  \item Additional recipes for react-navigation.
  \item A NativeWind/Tailwind theme providing the shadcn design tokens used by
        the \texttt{ui=shadcn} templates.
\end{itemize}

\textbf{Status:} \textbf{RECOMMENDED / POTENTIAL FUTURE WORK} --- explicitly not
committed plans.
